% Unit 081 English translation of chapter6.tex lines 1632--1739.
% Source: Methods of Algebra, Volume 2, commit
% 9a5803ff77dd3257484cb177f851a73770a59dd3 (CC BY 4.0).
% stable-unit-id: o014.aljabr2.chapter6.lhs-spectral-sequence-b
% Coverage: proof of the LHS theorem, the inflation--restriction sequence,
% Hopf's formula, and the multiplicative LHS structure; concludes Section 6.9.

% segment-id: o014.li2.chapter6.lhs-spectral-sequence-b.p001
\begin{proof}
	For~(i), it suffices to express \eqref{eqn:LHS-derived} concretely through
	the Grothendieck spectral sequence of
	Theorem~\ref{prop:Grothendieck-ss}.

	We now discuss~(ii). If $n=1$, the hypothesis is vacuous and the
	corresponding exact sequence is precisely the low-degree exact sequence in
	Theorem~\ref{prop:Grothendieck-ss}; its underlying principle is the edge
	computation in Remark~\ref{rem:edge-computing}. We next review and extend
	that argument to general $n$. The idea is to exploit the zero terms in the
	spectral sequence.

	It suffices to discuss the cohomological case. The hypothesis shows that
	the spectral sequence in~(i) satisfies
	\[ 1\leq q<n\implies
	E_2^{p,q}=E_3^{p,q}=\cdots=E_\infty^{p,q}=0. \]

	Recall that the direction of $d_r$ is
	$d_r^{p,q}:E_r^{p,q}\to E_r^{p+r,q-r+1}$ and that $E_{r+1}^{p,q}$ is the
	cohomology of $E_r^{p,q}$ with respect to $d_r$. Every $d$ leaving
	$E_{n+1}^{n+1,0}$ (or entering $E_{n+1}^{0,n}$) is zero. We therefore
	obtain the horizontal exact sequence in the diagram
	\[\begin{tikzcd}[row sep=small, column sep=scriptsize]
		& E_\infty^{0,n} \arrow[phantom,d,"\simeq" description,sloped]
		& & & E_\infty^{n+1,0}
		\arrow[phantom,d,"\simeq" description,sloped] & \\
		& \vdots \arrow[phantom,d,"\simeq" description,sloped]
		& & & \vdots \arrow[phantom,d,"\simeq" description,sloped] & \\
		0 \arrow[r] & E_{n+2}^{0,n} \arrow[r] &
		E_{n+1}^{0,n} \arrow[r,"{d^{0,n}_{n+1}}" inner sep=0.6em]
		\arrow[phantom,d,"\simeq" description,sloped] &
		E_{n+1}^{n+1,0} \arrow[r]
		\arrow[phantom,d,"\simeq" description,sloped] &
		E_{n+2}^{n+1,0} \arrow[r] & 0. \\
		& & \vdots \arrow[phantom,d,"\simeq" description,sloped] &
		\vdots \arrow[phantom,d,"\simeq" description,sloped] & & \\
		& \Hm^n(H,M)^{G/H} \arrow[r,"\sim"] & E_2^{0,n} & E_2^{n+1,0} &
		\Hm^{n+1}(G/H,M^H) \arrow[l,"\sim"'] &
	\end{tikzcd}\]
	The reason for the vertical isomorphisms is similar: once every $d$
	entering or leaving $E_r^{p,q}$ is zero, there is an isomorphism
	$E_r^{p,q}\simeq E_{r+1}^{p,q}$. In the same way one obtains
	\[ \Hm^n(G/H,M^H)\simeq E_2^{n,0}\simeq\cdots
	\simeq E_\infty^{n,0}. \]

	Next, recall that the target $\Hm^{p+q}(G,M)$ of the spectral sequence has
	a finite decreasing filtration
	$\cdots\supset\mathrm{F}^k\supset\mathrm{F}^{k+1}\supset\cdots$, with
	$E_\infty^{p,q}\simeq\gr^p\Hm^{p+q}(G,M)$. At the edge of the first
	quadrant, this gives
	\begin{align*}
		E_\infty^{n+1,0}&\simeq\mathrm{F}^{n+1}
		\subset\Hm^{n+1}(G,M),\\
		E_\infty^{0,n}&\simeq\Hm^n(G,M)/\mathrm{F}^1.
	\end{align*}

	How can one determine $\mathrm{F}^1\subset\Hm^n(G,M)$? On the line
	$p+q=n$, only $(n,0)$ and $(0,n)$ can possibly satisfy
	$E_\infty^{p,q}\neq0$. Thus the filtration on $\Hm^n(G,M)$ has the form
	\[ 0=\mathrm{F}^{n+1}
	\underbracket{\subset}_{E_\infty^{n,0}}\mathrm{F}^n
	=\cdots=\mathrm{F}^1
	\underbracket{\subset}_{E_\infty^{0,n}}\mathrm{F}^0
	=\Hm^n(G,M). \]
	Combining all the isomorphisms and exact sequences above gives the desired
	result.

	Assertion~(iii) is as expected, but the following argument is somewhat
	long and may be skipped. For brevity, we discuss only
	$\Hm^n(G,M)\to\Hm^n(H,M)^{G/H}$; the other cases are similar or easy.
	Recall the construction in \S\ref{sec:double-cplx-ss}. Take an injective
	resolution $M\to I$, with $I$ a complex of $G$-modules, and then take a
	Cartan--Eilenberg resolution $J$ of the complex of $G/H$-modules $I^H$
	(Theorem~\ref{prop:CE-resolution}). Regarding the complexes $I^G$ and so
	forth as double complexes concentrated on the horizontal axis gives a
	commutative diagram of double complexes of $\Bbbk$-modules
	\begin{equation}\label{eqn:LHS-SS-aux}\begin{tikzcd}
		& J^{G/H} \arrow[hookrightarrow,r] & J \\
		I^G \arrow[equal,r] & (I^H)^{G/H}
		\arrow[hookrightarrow,r] \arrow[u] & I^H. \arrow[u]
	\end{tikzcd}\end{equation}

	Following the convention of Proposition~\ref{prop:double-cplx-ss}, write
	$\Hm_{\mathrm{II}}\Hm_{\mathrm{I}}(\cdot)$ for the operation of first
	taking horizontal and then vertical cohomology of a double complex. Note
	the following points.
	\begin{itemize}
		\item Since the horizontal cohomology $\Hm_{\mathrm{I}}(J)$ of the
		Cartan--Eilenberg resolution gives an injective resolution of
		$\Hm^p(I^H)$ in column $p$ ($p\in\Z$), the morphism $I^H\to J$ induces
		the isomorphism
		\[ \Hm^n(H,M)=\Hm^n(I^H)
		=\Hm_{\mathrm{II}}\Hm_{\mathrm{I}}(I^H)^{n,0}
		\rightiso\Hm_{\mathrm{II}}\Hm_{\mathrm{I}}(J)^{n,0}. \]
		\item Remark~\ref{rem:CE-split} gives
		$\Hm_{\mathrm{I}}(J^{G/H})\simeq\Hm_{\mathrm{I}}(J)^{G/H}$. Since
		$\Hm_{\mathrm{I}}(J)$ vertically resolves each $\Hm^p(H,M)$, while
		$(\cdot)^{G/H}$ is left exact, applying
		$\Hm_{\mathrm{II}}\Hm_{\mathrm{I}}(\cdot)^{n,0}$ to
		$J^{G/H}\hookrightarrow J$ corresponds to
		\[ \text{the inclusion homomorphism}\qquad
		\Hm^n(H,M)^{G/H}\hookrightarrow\Hm^n(H,M). \]
		\item We have shown that $I^G\hookrightarrow I^H$ induces on $\Hm^n$
		(or, equivalently, on
		$\Hm_{\mathrm{II}}\Hm_{\mathrm{I}}(\cdot)^{n,0}$) the canonical
		homomorphism $\Hm^n(G,M)\to\Hm^n(H,M)$.
	\end{itemize}

	Filter the first-quadrant double complex $J^{G/H}$ decreasingly by the
	vertical coordinate (the $p$th step is the part with vertical coordinate
	$\geq p$). The corresponding spectral sequence is
	$(E_r^{p,q})_{p,q}$, where $r=0,1,\ldots$. Write $\Ker(d^n)_{n,0}$ for
	the projection of $\Ker(d^n)$ in the total complex $\tot(J^{G/H})$ onto
	the component $(n,0)$. The description of $Z_r^{p,q}$ in
	Proposition~\ref{prop:filtered-cplx-ss}, together with the observations
	above, gives
	\begin{multline}\label{eqn:LHS-SS-aux1}
		\Ker(d^n)_{n,0}\simeq Z_\infty^{0,n}
		\subset\cdots\subset Z_2^{0,n}\\
		\twoheadrightarrow E_2^{0,n}
		=\Hm_{\mathrm{II}}\Hm_{\mathrm{I}}(J^{G/H})^{n,0}
		\simeq\Hm^n(H,M)^{G/H}.
	\end{multline}

	We now describe $\Hm^n(G,M)\to\Hm^n(H,M)^{G/H}$. Choose a representative
	of an element of $\Hm^n(G,M)$ in
	$\Ker[I^{n,G}\to I^{n+1,G}]$, map it through $I^G\to J^{G/H}$ into
	$\Ker(d^n)_{n,0}$, and then follow \eqref{eqn:LHS-SS-aux1} to
	$\Hm_{\mathrm{II}}\Hm_{\mathrm{I}}(J^{G/H})^{n,0}$. To take its image
	further in $\Hm^n(H,M)$, apply
	$\Hm_{\mathrm{II}}\Hm_{\mathrm{I}}(\cdot)^{n,0}$ to the whole diagram
	\eqref{eqn:LHS-SS-aux} and compose along
	\begin{tikzpicture}[scale=0.5, baseline=(O)]
		\draw[->] (-1,0)--(0,0)--(0,1)--(1,1);
		\coordinate (O) at (0,0.2);
	\end{tikzpicture}.
	If instead we follow
	\begin{tikzpicture}[scale=0.5, baseline=(O)]
		\draw[->] (-1,0)--(1,0)--(1,1);
		\coordinate (O) at (0,0.2);
	\end{tikzpicture}
	and use the three observations above, the result is the canonical
	homomorphism $\Hm^n(G,M)\to\Hm^n(H,M)$. This proves~(iii).
\end{proof}

% segment-id: o014.li2.chapter6.lhs-spectral-sequence-b.p002
Remark~\ref{rem:LHS-SS-mult} will give a canonical and perhaps more natural
construction of the Lyndon--Hochschild--Serre spectral sequence.

% segment-id: o014.li2.chapter6.lhs-spectral-sequence-b.p003
\index{inflation--restriction exact sequence}
The cohomological exact sequence above is also called the
\emph{inflation--restriction exact sequence}, because the canonical
homomorphism $\Hm^n(G/H,M^H)\to\Hm^n(G,M)$ is commonly called the inflation
homomorphism, while $\Hm^n(G,M)\to\Hm^n(H,M)^{G/H}$ is called the restriction
homomorphism. The more complicated morphisms are
\[ \Hm^n(H,M)^{G/H}\to\Hm^{n+1}(G/H,M^H),\qquad
	\Hm_{n+1}(G/H,M_H)\to\Hm_n(H,M)_{G/H}. \]
They are both called \emph{transgression} maps, in keeping with the
transgression maps in the theory of characteristic classes.
\index{transgression}

% segment-id: o014.li2.chapter6.lhs-spectral-sequence-b.p004
As an application, we briefly derive a formula of Hopf for computing
$\Hm_2(G,\Z)$ from a presentation of the group $G$. Here we take $\Bbbk=\Z$
and give $\Z$ the trivial $G$-module structure.

% segment-id: o014.li2.chapter6.lhs-spectral-sequence-b.t001
\begin{corollary}[H.\ Hopf]\label{prop:Hopf-H2}
	\index[sym1]{[,]@$[\cdot,\cdot]$}
	Let $G$ be a group and $G\simeq F/R$, where $F$ is a free group and
	$R\lhd F$. Write $[F,R]$ for the subgroup of $R$ generated by the
	elements $[f,r]:=frf^{-1}r^{-1}$, with $f\in F$ and $r\in R$. This
	subgroup is normal in $R$ and
	$[F,R]\subset[F,F]:=F_{\mathrm{der}}$. There is an isomorphism
	\[ \Hm_2(G,\Z)\simeq(R\cap[F,F])/[F,R]. \]
\end{corollary}
\begin{proof}
	Proposition~\ref{prop:free-group-coh} gives $\Hm_2(F,\Z)=0$. Thus the
	exact sequence in Theorem~\ref{prop:LHS-SS} for $n=1$ becomes
	\[ 0\to\Hm_2(G,\Z)\to\Hm_1(R,\Z)_G
	\to\Hm_1(F,\Z)\to\Hm_1(G,\Z)\to0. \]
	Interpreting $\Hm_1$ through Example~\ref{eg:group-H1-further}, we obtain
	\[ 0\to\Hm_2(G,\Z)\to\left(\frac{R}{[R,R]}\right)_G
	\to\frac{F}{[F,F]}\to\frac{G}{[G,G]}\to0. \]

	The action of $G$ on $R/[R,R]$ is as follows. Choose a preimage $f\in F$
	of $g\in G$; one can verify that the action of $g$ comes from the
	conjugation action of $f$ on $R$ (exercise). For an arbitrary
	$\bar r\in R/[R,R]$, choose a preimage $r\in R$. Then
	$g\bar r-\bar r$ (with the group operation written additively) is the
	image of
	\[ frf^{-1}\cdot r^{-1}\in R
	\quad\text{(with the group operation written multiplicatively)}. \]
	As $f$ and $r$ vary, these elements generate $[F,R]$. Hence
	$(R/[R,R])_G\simeq R/[F,R]$; the rest is immediate.
\end{proof}

% segment-id: o014.li2.chapter6.lhs-spectral-sequence-b.r001
\begin{remark}[Multiplicative structure of the
Lyndon--Hochschild--Serre spectral sequence]\label{rem:LHS-SS-mult}
	In the cohomological case, write the spectral sequence of a $G$-module
	$M$ as $E_r^{p,q}(M)$. The cup-product operation to be introduced in
	\S\ref{sec:cup-product}, applied to the cohomology of $G/H$ and $H$,
	equips the $E_2$ term with a multiplicative structure
	\[ \cup:E_2^{p_1,q_1}(M_1)\otimes E_2^{p_2,q_2}(M_2)
	\to E_2^{p_1+p_2,q_1+q_2}(M_1\otimes M_2), \]
	while the convergence target of the spectral sequence also has a cup
	product, applied to the group $G$:
	\[ \cup:\Hm^{n_1}(G,M_1)\otimes\Hm^{n_2}(G,M_2)
	\to\Hm^{n_1+n_2}(G,M_1\otimes M_2). \]
	The question is whether these structures can be defined on every page and
	whether $d_r$ has good properties with respect to them, such as the
	Leibniz rule
	\[ d_r(\alpha_1\cup\alpha_2)
	=d_r(\alpha_1)\cup\alpha_2
	+(-1)^{p_1+q_1}\alpha_1\cup(d_r\alpha_2),
	\qquad \alpha_i\in E_r^{p_i,q_i}(M_i). \]

	In their original work \cite[pp.~118--119]{HS53}, Hochschild and Serre
	wrote down directly a canonical filtration on the normalized standard
	complex $\overline{C}(G,M)$ that is compatible with the cup product
	defined at the level of $\overline{C}(G,M)$. Consequently, every page
	inherits a multiplication with the properties above. Concretely, they use
	the decreasing filtration
	$\overline{C}(G,M)=\mathrm{F}^0\supset\mathrm{F}^1\supset\cdots$. If
	$p>n$, the degree-$n$ term of the complex $\mathrm{F}^p$ is $0$; otherwise,
	the degree-$n$ term of $\mathrm{F}^p$ is
	\[ \left\{\begin{array}{r|l}
		f\in\overline{C}^n(G,M) &
		(g_1,\ldots,g_n)\ \text{has at least }n-p+1\text{ entries in }H\\
		&\implies f(g_1,\ldots,g_n)=0
	\end{array}\right\}. \]
	The details are left as an exercise for this chapter; interested readers,
	or those who need them, may consult the original work.

	In particular, $\overline{C}(G,\Bbbk)$ becomes a filtered differential
	graded algebra, and the spectral sequence $E_r^{p,q}(\Bbbk)$ has the
	corresponding multiplicative structure (Definition~\ref{def:mult-ss} and
	Proposition~\ref{prop:dga-mult-ss}). From the topological viewpoint, this
	is entirely unsurprising.
\end{remark}
